% ── conclusion.tex · §6 · STATUS: skeleton (expand) ─────────────────
% Summary + future work + acknowledgements (ICSPA / CLEARSY).
\section{Conclusion and Future Work}\label{sec:conclusion}

We have presented \tool, which independently re-checks \ab's Predicate
Prover by reconstructing its proofs in \lp. Unlike approaches that
postulate each rule as an axiom, every \pp{} rule is encoded as a symbol
whose body \emph{proves} the rule sound against a small \bbook{} kernel,
so a reconstructed proof type-checks against a trusted base consisting of
that kernel and a handful of enumerated, argument-level solver
side-conditions. \tool{} covers nearly all of \pp's rule set and
reconstructs the non-arithmetic fragment of \ab's own proof corpus in
full.

The main avenue for future work is closing the arithmetic gap---either by
proving \pp's normalisation rules from an arithmetic theory on
\(\El\,\iota\) or by enriching the trace so the solver's reasoning becomes
reconstructible---and extending the parser to \pp's remaining argument
forms so the whole \texttt{prv} corpus can be attempted. Beyond coverage,
the soundness-proved encoding makes \pp{} a first-class citizen of the
\dedukti{} interoperability ecosystem, opening the door to exporting
\ab{} proofs to other proof assistants.

\subsubsection*{Acknowledgements.}
This work is supported by the ANR project ICSPA
(ANR-21-CE25-0015) in collaboration with CLEARSY.
% TODO: confirm grant number and any additional acknowledgements.
